% ===========================
% Conclusion
% ===========================

\vspace*{1cm}

\begin{center}
    \Large\textbf{Conclusion}
\end{center}

Pour conclure ce rapport, je souhaiterais revenir sur quelques difficultés et nouveautés que j'ai pu rencontrer au cours de ce stage.

La résolution par la méthode des caractéristiques n'a pas été simple. Malgré ma familiarité avec cette méthode, l'ajout d'une condition au bord constituait une nouveauté. La transformée de Laplace était également un outil complètement nouveau pour moi, même si sa prise en main n'a pas été particulièrement difficile en raison de sa proximité avec la transformée de Fourier.

L'étude par la théorie des semi-groupes a d'abord commencé par un travail de bibliographie. Une difficulté était de trouver le bon équilibre entre réaliser les démonstrations afin de bien comprendre les résultats et ne pas se noyer dans des raisonnements parfois très techniques dès la première lecture. La définition du cadre théorique pour le modèle de McKendrick-Lotka semblait simple au premier abord, mais a nécessité un second passage, notamment en raison d'un mauvais domaine de définition de l'opérateur en posant
\begin{equation*}
    D(A) := \Big\{ u \in W^{1,1}(0,a_{\max}) \ \ | \ \ u(0) = \int_0^{+\infty}\beta(a)u(a)\d a \Big\}
\end{equation*}
ce qui a posé problème lors de l'étude des propriétés du semi-groupe. La preuve de la densité n'a pas été simple non plus, car je ne suis pas très à l'aise avec la construction explicite de suites. Il en a été de même pour le contre-exemple sur la question du cadre général $W^{1,1}$.

L'étude du comportement asymptotique à l'aide de la théorie spectrale a, elle aussi, constitué une source de difficultés. Là encore, il ne fallait pas se perdre dans la bibliographie tant les nouveautés étaient nombreuses. De nombreuses propriétés devaient être démontrées et plusieurs éléments identifiés avant de pouvoir aboutir au résultat recherché. Il m'est souvent arrivé d'avoir l'impression de disposer de tous les ingrédients nécessaires sans pour autant savoir comment les assembler.

Enfin, je souhaitais terminer ce stage par l'étude d'un modèle non linéaire, d'où l'introduction d'un modèle de McKendrick-Lotka non linéaire basé sur le modèle EDO de Verhulst. L'étude de l'opérateur n'a pas posé de difficulté particulière puisqu'elle reposait sur de nombreux résultats déjà établis dans le cas linéaire. En revanche, l'étude des points d'équilibre et de leur stabilité constituait une véritable nouveauté. Ces notions n'avaient été abordées que très brièvement en troisième année de licence dans le cadre des EDO, et les raisonnements me paraissaient parfois peu intuitifs. Concernant la stabilité, il serait intéressant de poursuivre cette étude à l'aide des fonctions de Lyapunov. Une piste consisterait à rechercher une fonction de Lyapunov pour le point d'équilibre $K$ du modèle de Verhulst puis, si une telle fonction existe, tenter de l'adapter au modèle EDP afin de lever l'indétermination qui subsiste dans certains cas.

Mon objectif lors de ce stage était d'approfondir les notions d'analyse fonctionnelle abordées au cours de cette première année de master. De ce point de vue, ce stage a largement dépassé mes attentes. J'ai pu réinvestir de nombreuses notions étudiées cette année tout en découvrant de nouveaux outils mathématiques. L'étude du modèle de McKendrick-Lotka m'a notamment permis de découvrir la théorie des semi-groupes et la théorie spectrale, qui ont occupé une grande partie du temps consacré à ce travail et que j'ai particulièrement appréciées.

Ce stage m'a ainsi permis d'acquérir de nouvelles connaissances théoriques tout en développant ma manière d'aborder un problème de recherche, depuis l'étude bibliographique jusqu'à la mise en œuvre des outils mathématiques nécessaires à son analyse. Il constitue une expérience particulièrement enrichissante et une excellente préparation à la poursuite de mes études.